% ============================================================
% PAGE 2 — Related Work: Traditional Computer Vision Methods
% ============================================================

\section{Related Work}
\label{sec:related}

The history of object detection divides cleanly into two eras:
the \textit{classical} era characterized by hand-crafted
feature representations evaluated within a sliding-window
paradigm, and the \textit{deep-learning} era in which feature
extraction and classification are learned jointly from data.
This section reviews both, with emphasis on the design
decisions that created or alleviated the speed--precision
tension.

\subsection{Classical Feature-Engineering Detectors}

\subsubsection{Viola--Jones and Haar Cascades}

The Viola--Jones detector~\cite{viola2001rapid} introduced
three innovations that collectively enabled real-time face
detection for the first time.  The integral image representation
allowed rectangular Haar-like features to be computed in $O(1)$
time regardless of size, transforming an otherwise
computationally prohibitive feature extraction step into a
constant-time lookup.  The AdaBoost learning algorithm
selected a discriminative subset of features from a pool of
over 160\,000 candidates, yielding a compact but highly
effective classifier.  Most importantly, a cascade architecture
organized these classifiers from least to most expensive,
rejecting non-face windows at early stages before committing to
expensive later evaluations.  The cascade structure established
a fundamental principle that recurs throughout detection
history: \textit{coarse-to-fine filtering}, wherein cheap but
approximate tests gate access to more expensive but accurate
ones.

Despite its real-time throughput, Viola--Jones was narrowly
specialized: it required frontal, upright face images and
degraded sharply under occlusion, viewpoint change, or
cluttered backgrounds.  Attempts to generalize the method to
multi-class detection required training separate cascades for
each category, and accuracy on articulated objects remained
poor.

\subsubsection{Histogram of Oriented Gradients}

Dalal and Triggs~\cite{dalal2005histograms} proposed the
Histogram of Oriented Gradients (HOG) descriptor as a
rotationally tolerant, illumination-invariant feature for
pedestrian detection.  A detection window is divided into
overlapping spatial cells; within each cell, gradient
magnitudes are accumulated into orientation bins to form a
local histogram.  Cells are grouped into blocks and
block-level histograms are L2-normalized to achieve contrast
invariance.  The resulting descriptor is evaluated at every
scale and position of a sliding window by a linear support
vector machine (SVM), producing a dense response map from
which detections are extracted by non-maximum suppression
(NMS).

HOG improved robustness to illumination and minor deformations
relative to Haar features, but the sliding-window exhaustive
search scheme imposed an $O(S^2)$ computational burden with
respect to image scale factor $S$, making multi-scale detection
expensive.  Furthermore, the fixed spatial pooling grid assumed
rigid object geometry, precluding accurate localization of
flexible or articulated objects.

\subsubsection{Deformable Part Models}

Felzenszwalb \textit{et al.}~\cite{felzenszwalb2010object}
combined HOG-based appearance models with explicit geometric
deformation models in the Deformable Part Model (DPM)
framework.  Each object class is represented by a root filter
covering the full object and a set of part filters covering
discriminative regions; a latent SVM jointly learns filter
weights and part displacement costs, allowing the model to
score configurations in which parts deviate from their nominal
positions.  DPM won the PASCAL VOC detection challenge from
2007 to 2009, demonstrating superior performance on
deformable categories such as people and bicycles.

DPM's mixture models could represent multiple viewing aspects,
partially addressing the viewpoint limitation of Viola--Jones
and HOG.  However, inference remained slow---several seconds
per image at competitive accuracy settings---because the
matching over all feasible part placements involved solving a
dynamic programming problem at each scale.  Moreover, the
hand-crafted HOG features provided a ceiling on representational
power that could not be raised without replacing the underlying
feature extractor.

\subsection{The Transition to Deep Learning}

The ImageNet Large Scale Visual Recognition
Challenge~\cite{krizhevsky2012imagenet} catalyzed the shift
from hand-crafted to learned representations.  AlexNet reduced
the top-5 classification error from $26.2\%$ to $15.3\%$---a
margin so large it could not be attributed to algorithmic
improvements alone---by demonstrating that deep CNNs with
rectified linear units, dropout regularization, and GPU-accelerated
training could learn richer, more transferable features than
any previously designed descriptor.  The implications for
detection were immediate: if CNN features dramatically
outperformed HOG for classification, the same representations
could be expected to raise the ceiling on detection accuracy.

The challenge was architectural.  A classifier consumes an
entire image and produces a single label; a detector must
produce a variable-length set of boxes and labels, operating
efficiently across widely varying object scales and densities.
The following decade produced a series of architectures
progressively reconciling these requirements, as detailed in
the next section.

\subsection{Anchor-Based and Anchor-Free Paradigms}

A further axis of taxonomy cuts across both the two-stage and
single-stage lineages: whether predictions are made relative
to predefined \textit{anchor} boxes or directly in a
\textit{anchor-free} manner.  Anchor-based methods---Faster
R-CNN, SSD, YOLOv2 through YOLOv5---define a grid of reference
boxes at each spatial location and predict offsets from these
references.  Anchors reduce the regression problem to predicting
small residuals, which eases learning, but introduce
hyperparameters (aspect ratios, scales, stride) that must be
tuned to each dataset and can misalign with the distribution of
ground-truth boxes.

Anchor-free methods---CornerNet~\cite{law2018cornernet},
CenterPoint, FCOS---predict objects by detecting keypoints or
directly regressing box parameters from each spatial location
without reference boxes.  They eliminate anchor hyperparameters
and avoid the label-assignment ambiguities that arise when
multiple anchors overlap a single ground-truth box.  YOLOX
\cite{ge2021yolox} demonstrated that an anchor-free
reformulation of the YOLO architecture, combined with a
decoupled head and a dynamic label assignment strategy
(SimOTA), achieves higher mAP than anchor-based counterparts at
comparable FPS, a finding that directly motivates the
architectural choices of FusionDet described in
Section~\ref{sec:method}.
